%%%%%
%
\section{S�lection dans un vecteur selon un second vecteur (\texttt{float})}

\itemstyle{Op�rateur~:} \textsf{select\_X\_selon\_Y}
\label{OP-select_X_selon_Y}
\index{select\_X\_selon\_Y}

%%%%%OBJET%%%%
\itemstyle{Objet~:}
S�lection des donn�es d'un vecteur selon des
crit�res de gamme appliqu�s � un second vecteur de type \texttt{float}.

%%%%%ARGUMENTS%%%%
\bigskip
%%%%%ENTREES%%%%
\begin{tabarg}
\multicolumn{3}{l}{\titretab {Entr�es}}\\
\hline
$X$ & Vecteur, \texttt{float} & Param�tre � s�lectionner (unit� quelconque)\\
$Y$ & Vecteur, \texttt{float} & Param�tre selon lequel s'effectue la s�lection (unit� quelconque)\\
$Crit\_{min}$ & Coeff., \texttt{float} & valeur mini de $Y$ pour laquelle $X$ est consid�r� valide\\
$Crit\_{max}$ & Coeff., \texttt{float} & valeur maxi de $Y$ pour laquelle $X$ est consid�r� valide\\
$Y_{inval}$ & Coeff., \texttt{float} & valeur de $Y$ qui permettra d'�liminer �galement des valeurs de $X$ (�galement pour le cas o� la "valeur manquante" est mal reconnue (attribut misval non renseign� dans $Y$)\\
$Inv\_val$ & Coeff., \texttt{float} & valeur de remplacement de $X$ lorsque $Y$ ne correspond pas aux crit�res, le plus souvent "valeur manquante", soit \$Arg1.misval\\
\end{tabarg}

\medskip
%%%%%SORTIES%%%%
\begin{tabarg}
\multicolumn{3}{l}{\titretab {Sortie}}\\
\hline
$X_{valid}$ & Vecteur, \texttt{float} & $X$ lorsque $Y$ correspond aux crit�res. $Inv\_val$ dans les autres cas. 
(m�me unit� que $X$).\\
\end{tabarg}


%%%%%FORMULE%%%%
\itemstyle{Formule}

%%select_X_selon_Y
\[
X_{valid} = 
\begin{cases}
X \text{ si } Y \ne NaN \text{ et } Y \ne Y_{inval} \text{ et } Crit\_min \le Y \le Crit\_max \\
Inv\_val \text{ sinon}\\
\end{cases}
\]



%%%%%COMMENTAIRE%%%%
\itemstyle{Commentaire~:}

  Dans le cas o� le vecteur-crit�re (�$Y$�) est de type \texttt{int}, 
utiliser l'op�rateur
\hyperref[OP-select_flag_data]{\operfmt{select\_flag\_data}}.
Cet op�rateur prend en compte les cas o� $X$ est 1D et 2D.

%%%%%AUTEUR%%%%
\itemstyle{Auteur~:}
CNRM/GMEI/TRAMM

%%%%%REFERENCE%%%%
\itemstyle{R�f�rences~:}

